
% ==========================================
% CHAPTER 8
% ==========================================
\chapter{Classes and Object Oriented Programming} % 

\section{Defining Classes} % 
Objects, or Classes, are useful for a variety of reasons; primarily we use them to easily organize groups of functions that act on a specific ``object'' that we define. %  
There are two main components to a defined class: attributes and methods. %  
Attributes are properties which are created when an instance of a class is initialized. %  
Methods are functions which each instance of a class carries and can utilize. %  

Say we want to create a class called \texttt{Planet} which has certain attributes such as planet name, revolution period, and mass. % 

\textbf{Example 8.1 Defining a Planet Class} % 
\begin{lstlisting}
class Planet(object):
    def __init__(self, planet_name, rev_period, mass):
        self.planet_name = planet_name
        self.rev_period = rev_period
        self.mass = mass
\end{lstlisting} % 

The first function in a class is always the init function (which has a double underscore before and after the word init). %  
The first argument of init is always ``self''. %  
Once we have created an \texttt{\_\_init\_\_} function, we can put other functions in our class as well. %  
These are known as methods. %  

\textbf{Example 8.2 Methods} % 
\begin{lstlisting}
class Planet(object):
    def __init__(self, planet_name, rev_period, mass):
        self.system_name = '2014-B178h'
        self.planet_name = planet_name
        self.rev_period = rev_period
        self.mass = mass

    def semimajoraxis(self):
        return self.rev_period**2 * (2./3)
\end{lstlisting} % 

\subsection{Subclasses} % 
It is possible to make subclasses of your own, which inherit all the attributes and methods of their parent class, while having some specialized methods and attributes of their own. %  
For example, we can make a subclass of \texttt{Planet} called \texttt{Dwarf}: % 

\textbf{Example 8.3 Subclasses} % 
\begin{lstlisting}
class Dwarf(Planet):
    def describe(self):
        return self.name + ':' + 'Mass=' + str(self.mass)
\end{lstlisting} % 

\section{Initializing an Instance of a Class} % 
We can initialize a variable as an object (say, a planet), like this: % 

\begin{lstlisting}
earth = Planet('Earth', 24, 5.9E24)
\end{lstlisting} % 

We can access those properties of any of our class objects using dot notation: % 

\begin{lstlisting}
print(earth.mass)
\end{lstlisting} % 

Now that we have the instance of the \texttt{Planet} class initialized as the variable \texttt{earth}, we can also execute the methods that we have defined for the given class. %  
For example, using the method we defined above, we can call: % 

\begin{lstlisting}
print(earth.semimajoraxis())
\end{lstlisting} % 
